% Conclusiones
\section{Conclusiones}

\begin{frame}
\frametitle{Hallazgos clave}
\begin{enumerate}
    \item El benchmark MCO/Probit alcanza AUC-ROC = 0.7232. Es un piso difícil de superar con sólo 4 \textit{features}.
    \item Los modelos ML por defecto \textit{no} superan al benchmark.
    \item El \textbf{GBM tuneado} sí lo hace: AUC = \textbf{0.7278}, F1 = \textbf{0.6129}.
    \item SHAP: edad + edad² $\approx 77\%$ de la importancia. Modelo \emph{honestamente delgado}.
    \item Robustez Y alternativa: relaciones estructurales estables.
\end{enumerate}
\end{frame}

\begin{frame}
\frametitle{Lecciones metodológicas}
\begin{itemize}
    \item La flexibilidad sin regularización introduce varianza que opaca la ganancia de sesgo.
    \item El \textit{tuning} disciplinado importa más que el algoritmo en datos chicos.
    \item La interpretabilidad no es opcional: define qué tan utilizable es el modelo para política.
    \item Documentar limitaciones (Y discrepante, submuestra con precio) es parte del trabajo, no un anexo.
\end{itemize}
\end{frame}

\begin{frame}
\frametitle{Próximos pasos}
\begin{itemize}
    \item Reconciliar la construcción de \texttt{consumo\_12m} con \texttt{K\_03} directo.
    \item Ampliar el \textit{feature set} con variables de entorno (geografía, red social, ingreso del hogar).
    \item Estimar elasticidad-precio con un diseño que corrija el sesgo de selección.
    \item Acoplar el modelo de propensión con proyecciones poblacionales del DANE para simulación tributaria.
\end{itemize}
\end{frame}

\begin{frame}[plain]
\centering
\Huge \textbf{Preguntas y discusión}
\vspace{0.6cm}
\\[0.4cm]
\normalsize
\textit{Tupaz, Quintero, Rodrigues — Universidad EAFIT}\\
\textit{Profesora supervisora: Paula María Almonacid Hurtado}
\end{frame}
